\chapter{Experimental Setup}

\section{Hardware and Cluster Environment}
The target environment is the macOS/OpenMPI cluster described by the repository hostfiles. GPU-oriented experiments use \code{configs/hosts\_macos\_gpu}; CPU/core-count experiments use \code{configs/hosts\_macos\_core} and weighted hostfiles when process counts exceed the three-rank GPU layout. Exact CPU, memory, process-slot, network, OpenMPI, and Python details must be taken from the final hardware evidence logs. They are intentionally not invented in this draft.

\begin{table}[H]
\centering
\caption{Hardware evidence to be filled from cluster logs}
\label{tab:hardware}
\begin{tabularx}{\textwidth}{L{2.4cm} L{2.2cm} L{2.6cm} L{2.0cm} X}
\toprule
Host & Role & CPU/SoC & Memory & Evidence source \\
\midrule
master & Pending log & Pending log & Pending log & final evidence logs \\
node1 & Pending log & Pending log & Pending log & final evidence logs \\
node2 & Pending log & Pending log & Pending log & final evidence logs \\
\bottomrule
\end{tabularx}
\end{table}

\section{Software Environment}
The implementation uses C++17, OpenMPI, a Python virtual environment, Ultralytics YOLO11n, and helper scripts under \path{scripts/runtime/}. The final report must record the actual OpenMPI version, Python version, YOLO package version, model path, detector device, and source video path from the run logs.

\section{Dataset and Input Size Definition}
The input size is the number of processed video frames \(N\). The input-size selection experiment runs several candidate frame counts and chooses an \(N\) whose wall-clock runtime is approximately 120--180 seconds on the target configuration. Speedup experiments then use \(2N\) frames to reduce startup noise.

\section{Metrics and Timing Methodology}
The report uses the CSV fields generated by the implementation:
\begin{itemize}
    \item \code{total\_ms\_with\_comm}: wall-clock runtime including MPI communication;
    \item \code{total\_ms\_without\_comm}: compute-oriented estimate excluding communication;
    \item \code{compute\_ms}, \code{yolo\_ms}, \code{comm\_ms}, \code{idle\_ms}: rank-level timing;
    \item \code{tasks\_done}: task count per rank;
    \item \code{correctness\_pass}: serial-versus-parallel verification status when \code{--verify 1}.
\end{itemize}
Speedup and efficiency are:
\begin{equation}
    S_p = \frac{T_1}{T_p}, \qquad E_p = \frac{S_p}{p}.
\end{equation}
Communication overhead is:
\begin{equation}
    \frac{T_{comm}}{T_{wall}}.
\end{equation}

\section{Required Experiment Matrix}
\begin{table}[H]
\centering
\caption{Required YOLO-only experiment matrix}
\label{tab:experiment-matrix}
\begin{tabularx}{\textwidth}{L{3.1cm} L{3.1cm} L{3.2cm} X}
\toprule
Experiment & Main setting & Required output & Report purpose \\
\midrule
Correctness & Static and dynamic, \code{--verify 1} & correctness CSVs & Verify parallel output against serial baseline. \\
Input size & Multiple \(N\), \(P=12\), dynamic & \code{raw/find\_N.csv} & Choose \(N\) near 120--180 s. \\
Granularity & Selected \(N\), one or more tile grids & \code{rank\_metrics.csv} & Per-rank task, compute, comm, idle analysis. \\
Scheduling & Same \(N,P\), static vs dynamic & scheduler CSV & Compare runtime and imbalance. \\
Speedup & \(2N\), \(P=1,2,4,8,12,24\) if possible & \code{raw/speedup.csv} & Compute \(S_p\) and \(E_p\). \\
Communication & Derived from summaries and rank metrics & communication CSV & Quantify \(T_{comm}/T_{wall}\). \\
\bottomrule
\end{tabularx}
\end{table}
